\textbf{11-1.} S'系相对 S 系以速度 $v = 0.8c$ 沿 $x$ 轴运动。S'系和 S 系的原点在 $t = t' = 0$ 时重合。一事件在 S'系中发生在 $x' = 300 \mathrm{~m} (y' = z' = 0)$ , $t' = 2 \times 10^{-7} \mathrm{~s}$ 。求该事件在 S 系中发生的空间位置 $x$ 和时间 $t$ 。

\vfill

\textbf{11-2.}在一惯性系的同一地点,先后发生两个事件,其时间间隔为 $0.2 \mathrm{~s}$ , 而在另一惯性系中测得此两事件的时间间隔为 $0.3 \mathrm{~s}$ , 求两惯性系之间的相对运动速率。

\vfill
\newpage

\textbf{11-3.}S'系相对 S 系以恒定速度沿 x 轴运动。在 S 系中两事件发生在同一时刻, 并沿 x 轴相距 $2 \times 10^{4}$ m。而在 S'系中, 测得两事件的空间间隔为 $3 \times 10^{4}$ m, 试问在 S'系中测得的两事件的时间间隔是多少?

\vfill

\textbf{11-4.}一火箭飞船经地球飞往某空间站,该空间站相对地球静止,与地球之间的距离为 $9.0 \times 10^{9} \mathrm{~m}$ 。地球上的钟和空间站上的钟是同步校准的。当火箭飞船飞经地球时,宇航员将飞船上的钟拨到与地球上的钟相同的示数。当火箭飞船飞经空间站时,宇航员发现飞船上的钟比空间站上的钟慢了 $3 \mathrm{~s}$ 。求火箭飞船的飞行速率。

\vfill
\newpage

\textbf{11-5.}一飞船以 $v = 0.6c$ 的速率沿平行于地面的轨道飞行。飞船上沿运动方向放置一根杆子，在地面上的人测得此杆子的长度为 $l$ ，求此杆子的本征长度 $l_{0}$ 。

\vfill

\textbf{11-6.}一道闪光从 $O$ 点发出, 在 $P$ 点被吸收。在 S 系中, $OP$ 具有长度 $l$ 且

同 x 轴成 $\theta$ 角, 如习题 11-6 图所示。在相对于 S 系以恒速 v 沿 x 轴运动的 S' 系中:

(1) 从光的发出到吸收相隔多长时间 $\Delta t'$ ?

(2) 光的发出点 O 到吸收点 P 的空间间隔 $l'$ 是多大?

\vfill
\newpage

\textbf{11-16.}两根相互平行的米尺,各以 $v = \frac{3}{5} c$ 的速率相向运动,运动方向平行于尺子。求任一尺子上的观察者测量另一尺子的长度。

\vfill


\textbf{11-17.}静长为 $L$ 的车厢, 以恒定速率 $v$ 沿地面向右运动。自车厢的左端 $A$ 点发出一光信号, 经右端 $B$ 点镜面反射后回到 $A$ 端。

(1) 在车厢里的人看来, 光信号经多少时间 $\Delta t_{1}^{\prime}$ 到达 $B$ 端? 从 $A$ 发出经 $B$ 反射后回至 $A$ , 共需多少时间 $\Delta t^{\prime}$ ?

(2) 在地面上的人看来, 光信号经多少时间 $\Delta t_{1}$ 到达 $B$ 端? 从 $A$ 发出经 $B$ 反射后回至 $A$ , 共需多少时间 $\Delta t$ ?

\vfill\newpage

\textbf{11-18.}一艘静长为 $90 \mathrm{~m}$ 的飞船以速度 $v = 0.8 c$ 飞行。当飞船的尾部经过地面上某信号站时, 该信号站发出一光信号。

(1) 当光信号到达飞船头部时, 飞船头部离地面信号站的距离为多远?

(2) 按地面上的时间, 信号从信号站发出共需多少时间 $\Delta t$ 才能到达飞船头部?

\vfill


\textbf{11-19.}三艘飞船 A、B、C 沿同一直线飞行。B 船上的观察者发现，A、C 两船以 $v = 0.7c$ 的相同速率远离 B 船。

(1) 求 A 船上的观察者观测到的 B、C 两船的速率；

(2) 若 B 船相对地面的飞行速率为 $v_{\mathrm{B}} = 0.9c$ , 求 A、C 两船相对地面的飞行速率。

\vfill\newpage

\textbf{11-20.}在实验室里测得一根沿 $x$ 方向运动的棒与 $x$ 轴的夹角 $\theta_{1} = 45^{\circ}$ 。在相对实验室参考系以 $v = 0.6c$ 的速度沿 $x$ 方向运动的另一参考系 $\mathbf{S}'$ 系中, 测得此夹角 $\theta' = 35^{\circ}$ 。

(1) 求棒相对实验室参考系的运动速度；

(2) 棒相对 S' 系的运动速度为多大?

(3) 在棒静止的参考系 S" 中, 棒与 x" 轴的夹角 $\theta'$ 为多大?

\vfill


\textbf{11-21.}一航天飞机沿 $x$ 方向飞行, 接静止的参考系中, 飞机的飞行速度为 $v$ , 恒星发出光信号的方向与飞机的轴成 $\theta$ 角, 如习题 11-21 图所示。

(1) 在飞机静止的参考系中此角 $\theta'$ 为多大?

(2) 若飞机的前端有一半球形的观察室, 飞船上的人能看到所有 $\theta' < \frac{\pi}{2}$ 的恒星。试证明: 若 $v \to c$ , 则飞机上的人几乎能看到所有的恒星。



\vfill\newpage

\textbf{11-22.}光在介质中的传播速率 $c_{n} = \frac{c}{n}$ , 这里 $n$ 是介质的折射率。设光在流动的水中传播, 水的流速为 $v$ , 折射率为 $n_{0}$ , 在水静止的参考系中, 光沿水流方向的传播速度为 $c'_{n} = \frac{c}{n_{0}}$ 。求在实验室参考系中测得的该光速 $c_{n}$ 。

\vfill


\textbf{11-23.}一艘火箭飞船以 $v = 0.8c$ 的速度飞经地球, 飞船和地球上的观察者一致同意该事件发生在中午 12:00。

(1) 按照火箭飞船上的时钟读数,该飞船于 12:30 飞经一个行星际宇航站,该站相对地球固定,其时钟指示地球时间。试问这一事件在该站什么时间发生?

(2) 在地球坐标上该站离地球多远?

(3) 在飞船时间 12:30(即飞船飞经该宇航站时), 飞船用无线电向地球发回报告。试问按地球时间, 地球何时收到信号?

(4) 如果地面站立即回答,试问按飞船时间,飞船何时接到答复?

\vfill\newpage

\textbf{11-24.}一宇航员乘坐宇宙飞船去星际航行。如果他从静止开始离开地球，在他的瞬间静止参考系（即与某一瞬间该飞船相对地球运动速度相同的惯性系，不同的时刻对应不同的惯性系）中持续以 $g = 9.8 \mathrm{~m} / \mathrm{s}^{2}$ 的加速度做加速运动。

(1) 在地球时间 t 时, 他已经飞过了多长的距离?

(2) 在他的速率达到 $\frac{1}{2} c$ 之前, 他已经飞行了多长时间?

\vfill


\textbf{11-25.}氢原子发出的一条光谱线 $\mathrm{H}_{\delta}$ 的波长 $\lambda = 0.4101 \times 10^{-6} \mathrm{~m}$ 。在极隧射线管中, 氢原子的速率可达 $v = 5 \times 10^{5} \mathrm{~m} / \mathrm{s}$ 。当氢原子以此速率向着观察者飞行时,此谱线的波长将增大还是减小?波长将改变多少?已知光速 $c=3\times10^{8}$ m/s。

\vfill\newpage

\textbf{11-26.}对遥远星系发来的光所作的光谱分析表明, 光的谱线非常明显地移向可见光谱的红端, 这可解释作为光源的这些星系正向远离地球的方向退行而引起的多普勒效应, 故称退行红移。比如钾光谱中有一对容易辨认的吸收线 (K 线和 H 线), 其谱线的波长在 $395.0 \mathrm{~nm}$ 附近。而来自牧夫星座某个星云的光中, 我们在波长为 $447.0 \mathrm{~nm}$ 处发现了这两条谱线。试求该星云的退行速度。

\vfill


\textbf{11-27.}当一光源以速率 $v_{1}$ 向地球靠近时, 在地球静止系 S 中的人 A 看到光源发出的是绿光 ( $\lambda_{1} = 500.00 \mathrm{~nm}$ ), 而在相对地球以速率 $v_{2}$ 运动 (与光源的运动沿同一直线) 的参考系 S' 中的人 B 看到的却是红光 ( $\lambda_{2} = 600.00 \mathrm{~nm}$ )。当该光源以相同的速率 $v_{1}$ 远离地球时, A 看到的是红光 ( $\lambda_{2} = 600.00 \mathrm{~nm}$ )。

(1) 求 $v_{1}$ 、 $v_{2}$ 的值；

(2) 当光源以 $v_{1}$ 远离地球时, B 看到光的波长为多大?

\vfill\newpage

\textbf{11-28.}半人马座 $\alpha$ 星与地球相距 $4.3\ l.y.$ (光年), 两个孪生兄弟中的一个 A 乘坐速度为 $0.8 c$ 的宇宙飞船去该星旅行。他在往程和返程途中每隔 $0.01 a$ 的时间 (飞船静止参考系的时间) 发出一个无线电信号。另一个留在地球上的孪生兄弟 B 也在相应过程中每隔 $0.01 a$ 的时间 (地球静止参考系的时间) 发出一个无线电信号。

(1) 在 A 到达该星以前, B 收到多少个 A 发出的信号?

(2) 在 A 到达该星以前, A 收到多少个 B 发出的信号?

(3) A 和 B 各自共收到多少个从对方发出的信号?

(4) 当 A 返回地球时, A 比 B 年轻了几岁? 试证明两孪生兄弟都同意此观点。

\vfill


\textbf{11-29.}一粒子的动能等于静能的一半,试求其运动速度。

\vfill\newpage

\textbf{11-30.}一沿 $x$ 轴正方向运动的光子具有 $400 \mathrm{MeV}$ 的能量, 另一沿 $y$ 轴正方向运动的光子具有 $300 \mathrm{MeV}$ 的能量, 若有一粒子的总能量和动量与此两光子的总能量和总动量相同。试求:

(1) 该粒子的静质量 $m_{0}$ ;

(2) 该粒子的运动速度 v。

\vfill


\textbf{11-31.}在 S 参考系中观测得一粒子的总能量为 500 MeV, 动量为 400 MeV/c。而在 S' 参考系中观测得该粒子的总能量为 583 MeV。试求:

(1) 该粒子的静能；

(2) 在 S' 系中观测到的该粒子的动量；

(3) $S'$ 系相对 $S$ 系的运动速度。

\vfill\newpage

\textbf{11-32.}两静质量相同的粒子,一个处于静止状态,另一个的总能量为其静能的4倍。两粒子发生碰撞后黏合在一起,成为一复合粒子。求此复合粒子的静质量与碰撞前单个粒子的静质量之比。

\vfill


\textbf{11-33.}如习题 11-33 图所示, 一个以 $0.8c$ 的速率沿 $x$ 轴正方向运动的粒

子衰变成两个静质量均为 $m_{0}$ 的粒子, 其中一个粒子以 0.6c 的速率沿 -y 方向运动。设衰变前粒子的静质量为 $m'_{0}$ 。试求

(1) 另一个粒子的运动速率 u 和方向（用图中的 $\theta$ 表示）；

(2) $\dfrac{m_{0}}{m'_{0}}$ 的值。



\vfill\newpage

\textbf{11-34.}动能为 $6 \times 10^{3} \mathrm{MeV}$ 的质子与静止质子碰撞时, 能形成质子-反质子对。若用动能相同的两质子对撞来实现此反应, 试求其动能的最小值。已知质子的静能为 $938 \mathrm{MeV}$ 。

\vfill


\textbf{11-35.}设有一处于激发态的原子以速率 $v$ 运动。当其发射一能量为 $E'$ 的光子后衰变至其基态, 并使原子处于静止状态, 此时原子的静质量为 $m_0$ 。若激发态比基态能量高 $E_0$ , 试证明:

$$
E ^ {\prime} = E _ {0} \left(1 + \frac {E _ {0}}{2 m _ {0} c ^ {2}}\right)
$$

\vfill\newpage

\textbf{11-36.}太空火箭(包括燃料)的初始质量为 $m'_{0}$ , 从静止起飞, 向后喷出的气体相对火箭的速度 $u$ 为常量。当火箭相对地球速度为 $v$ 时, 其静止质量为 $m_{0}$ , 试求 $m_{0}/m'_{0}$ 的值与速度 $v$ 之间的关系。设此过程中忽略其他天体的引力影响。

\vfill


\textbf{11-37.}光子火箭从地球起飞时初始静质量(包括燃料)为 $m'_{0}$ , 向相距为 $d = 1.8 \times 10^{6} \mathrm{l} \cdot \mathrm{y}$ . 的仙女座星云飞行。要求火箭在火箭时间 $25 \mathrm{~a}$ 后到达目的地。设不计其他星球的引力影响。

(1) 忽略火箭加速和减速所需的时间, 求火箭所需的速度;

(2) 设到达目的地时火箭的静质量为 $m_{0}$ , 求 $m^{\prime}_{0} / m_{0}$ 的最小值。

\vfill\newpage

\textbf{11-38.}一个 $\alpha$ 粒子以速率 $v_{1} = \frac{4}{5} c$ 进入厚度 $d = 0.35 \mathrm{~m}$ 的水泥防护墙, 当此粒子从墙的另一面穿出时, 速率减小为 $v_{2} = \frac{5}{13} c$ 。已知 $\alpha$ 粒子的静质量 $m_{0} = \frac{2}{3} \times 10^{-26} \mathrm{~kg}$ 。设墙对粒子的作用力是常量。

(1) 在墙静止的参考系 S 中作用力 $F_{0}$ 的值为多大?

(2) 在以速率 $v_{1}$ 沿粒子运动方向相对墙运动的参考系 $\mathbf{S}'$ 中作用力 $F'_{0}$ 的值为多大?

(3) 在 S 与 S' 系中粒子穿过墙各需多长时间?


\vfill
\newpage